%%
%% SECTION: Qualitative Analysis of Generated Responses Under SVD Truncation
%%
%% Drop this block into neurIPS_paper_final.tex before \section{Discussion}
%%

\section{Qualitative Analysis of Generated Responses Under SVD Truncation}

While KL divergence and effective rank quantify \textit{how much} the output distribution shifts,
they do not reveal \textit{what changes} in the generated text.
In this section, we perform a systematic qualitative and quantitative analysis of the 42 generated
responses under rank-$k$ truncation. We measure Jaccard word overlap, shared prefix length, step-count
preservation, and word-count drift, and we identify five distinct generation archetypes that collectively
explain the W-shape divergence pattern and the role of secondary singular modes.

\subsection{Aggregate Text Metrics}

Table~\ref{tab:text_metrics} summarises the four text-level metrics at $k=1$ for all six BTL levels.

\begin{table}[h]
\centering
\caption{Text divergence metrics at $k=1$ (mean over 7 prompts per level). Jaccard = word-bag overlap;
Prefix\% = fraction of generated tokens identical before first divergence;
$\Delta$W = word-count change vs.\ baseline; Step\% = fraction of prompts preserving step count.}
\label{tab:text_metrics}
\begin{tabular}{lcccc}
\toprule
\textbf{BTL Level} & \textbf{Jaccard} & \textbf{Prefix \%} & \textbf{$\Delta$W} & \textbf{Step\%} \\
\midrule
Remembering   & \textbf{0.658} & \textbf{32.0\%} &  $+0.3$ & \textbf{71\%} \\
Analyzing     & 0.637          & 31.7\%          &  $-1.9$ & 57\%          \\
Applying      & \textbf{0.666} & 19.5\%          &  $-0.6$ & 29\%          \\
Understanding & 0.508          & 21.0\%          &  $-0.6$ & 43\%          \\
Creating      & 0.506          & 19.5\%          &  $-5.1$ & 43\%          \\
Evaluating    & 0.366          & 17.3\%          & $+7.7$  & 29\%          \\
\bottomrule
\end{tabular}
\end{table}

Three observations stand out immediately.
\textbf{(i)} Applying achieves the highest Jaccard similarity (0.666) despite a low prefix match (19.5\%)
and low step-structure preservation (29\%), a dissociation explained in Section~\ref{sec:applying}.
\textbf{(ii)} Evaluating is the only level where truncation \textit{inflates} response length ($+$7.7 words),
for reasons explained in Section~\ref{sec:evaluating}.
\textbf{(iii)} The overall word count is remarkably stable across levels (all $|\Delta W| < 8$),
confirming that SVD truncation does not suppress generation fluency.

%% ── ── ── ── ── ── ── ── ── ── ── ── ── ── ── ── ── ── ──
\subsection{Remembering: Template Preservation with Metadata Compression}

Rote recall prompts (e.g., \textit{``List the US Presidents who served during the 20th century''})
produce the most stable outputs under truncation.
At $k=1$, the full step-by-step scaffold is preserved in 71\% of cases---the highest of any level.
What does change is the compression of relative-clause qualifiers:

\begin{quote}
\small
\textbf{Baseline:} ``Recall the next US President \underline{who served after the first one}. Step 4: Remember the name...'' \\
\textbf{$k=1$:} ``Recall the next US President. Step 4: Remember the name...''
\end{quote}

The truncated model drops the subordinate clause (``who served after the first one'') while
preserving the action verb, the entity, and the ordering structure.
This is not comprehension failure---the sentence remains grammatically and factually correct.
It is \textit{syntactic metadata compression}: the primary singular mode ($v_1$) encodes the
high-frequency action-entity pairs (Recall, President), while secondary modes ($v_2,\,v_3$)
encode relational specificity.
Remembering tasks are therefore naturally \emph{rank-1 compatible}: the information geometry
of factual retrieval is intrinsically low-dimensional.

%% ── ── ── ── ── ── ── ── ── ── ── ── ── ── ── ── ── ── ──
\subsection{Understanding: Declarative-to-Temporal Framework Shift}
\label{sec:understanding}

Understanding-level prompts (e.g., \textit{``Explain how a four-stroke internal combustion engine works''})
show the most dramatic structural reorganisation.
The same factual content is preserved---intake stroke, piston moving downward, vacuum formation---but
the discourse framing shifts entirely across truncation levels:

\begin{quote}
\small
\textbf{Baseline:} \textit{``The first stroke \underline{is called} the intake stroke. During this stroke, the piston...''}\\
\textbf{$k=1$:}  \textit{``The engine \underline{starts with} the intake stroke. The piston...''}\\
\textbf{$k=5$:}  \textit{``The \underline{first step in} a four-stroke engine is the intake stroke. During this stroke...''}
\end{quote}

The baseline adopts a \textbf{definitional register} (``is called''), a formal taxonomic style.
The $k=1$ output shifts to a \textbf{temporal narrative register} (``starts with''),
and $k=5$ falls between them with an instructional framing (``first step in'').
Critically, prefix match collapses to 21\% and step-structure preservation to 43\%.

This finding explains why Understanding carries the \textbf{highest KL divergence} ($D_\mathrm{KL}=0.283$)
despite being a mid-range cognitive level.
Sustained explanatory discourse requires the model to maintain a \emph{consistent rhetorical register} over
many generation steps---``is called'' at step 1 conditions ``is called'' at steps 2, 3, 4 via in-context
imitation.
When $v_2\dots v_4$ are removed, the model cannot sustain this register.
It falls back to the highest-frequency discourse pattern in pre-training: sequential temporal narration.
Register coherence across multiple steps is, therefore, a \emph{secondary-mode phenomenon}.


%% ── ── ── ── ── ── ── ── ── ── ── ── ── ── ── ── ── ── ──
\subsection{Applying: Formulaic Immunity and the Rank-1 Floor}
\label{sec:applying}

Mathematical and procedural prompts (e.g., \textit{``Calculate the trajectory of a projectile at 45$^\circ$,
50\,m/s''}) exhibit a qualitatively different behaviour: \textbf{perfect output identity at $k=1$}---verified
token-for-token across the full 200-token generation.

The mechanism is straightforward: established physical formulae create near-deterministic next-token
distributions.
When the model generates ``Vx = V\,$\cdot$\,cos'', the only valid continuation within the kinematics
vocabulary is ``(theta)''.
There is essentially no probability mass for secondary tokens; the attention distribution collapses
to a rank-1 manifold \textit{by the nature of the domain itself}.

This defines what we term the \textbf{rank-1 floor}: a task-induced lower bound on attentional
dimensionality below which any truncation is information-lossless.
Mathematical syntax, physical constants, chemical equations, and rigid procedural protocols all reside
at or near this floor.

The apparent paradox in Table~\ref{tab:text_metrics} is resolved: Applying achieves high Jaccard (0.666)
because domain-specific vocabulary (cos, sin, $\theta$, m/s) is reproduced identically across all models,
but low step-structure preservation (29\%) because the \textit{scaffolding} around calculations---
how many sub-steps to enumerate, whether to use bullet points or inline notation---depends on
secondary singular modes.
Knowledge of what to compute is rank-1; knowledge of how to present it is not.

%% ── ── ── ── ── ── ── ── ── ── ── ── ── ── ── ── ── ── ──
\subsection{Analyzing: Lexical Substitution with Preserved Logic Chain}

Comparative analytical prompts (e.g., \textit{``Analyze the thematic differences between Marvel and DC''})
show strong robustness, with Jaccard=0.637 and prefix=31.7\%---closely matching Remembering.

The truncation effect is limited to fine-grained lexical synonymy:

\begin{quote}
\small
\textbf{Baseline:} ``...known for their \underline{darker, more mature} tone...''\\
\textbf{$k=1$:}  ``...known for their \underline{darker and more mature} tone...''\\
\textbf{$k=5$:}  ``...often have a \underline{darker and more serious} tone...''
\end{quote}

The binary comparative frame (Marvel vs.\ DC, tone vs.\ world-building) creates
predictable high-probability anchor positions in the generation trajectory.
Secondary modes encode slight synonym preferences (``mature'' vs.\ ``serious'',
``often focusing'' vs.\ ``focusing'') but do not alter the logical argument structure.
Notably, $k=5$ produces the most deviant word choices---``more serious''
instead of the more frequent ``more mature''---suggesting that \emph{higher-rank
reconstructions can sometimes select rarer but equally valid alternatives}
as the distributional mass spreads across more modes.

%% ── ── ── ── ── ── ── ── ── ── ── ── ── ── ── ── ── ── ──
\subsection{Evaluating: Interrogative Collapse and Modal Inversion}
\label{sec:evaluating}

Evaluative prompts produce the most consequential structural shift in the entire dataset.
Consider the four-day workweek prompt:

\begin{quote}
\small
\textbf{Baseline:} ``\underline{Consider} the economic impact. Step 2: \underline{Analyze} the psychological benefits.
Step 3: \underline{Evaluate} the logistical challenges. Step 4: \underline{Weigh} the pros and cons...'' \\[4pt]
\textbf{$k=1$:}  ``The four-day workweek is a proposal that aims to increase productivity...
Step 2: To evaluate its effectiveness, we need to consider the economic impact.
\underline{Will it lead to increased productivity?} Step 3: We also need to consider the psychological impact.
\underline{Will employees be more motivated?}...''
\end{quote}

The baseline uses \textbf{imperative command mode}: a pre-planned evaluation agenda executed as
concise directives (``Consider...'', ``Analyze...'', ``Evaluate...''). 
The $k=1$ output undergoes \textbf{complete modal inversion into interrogative framing}:
the model no longer \textit{executes} an evaluation plan but instead \textit{formulates the questions}
that such a plan would answer.

This distinction is not superficial.
Imperative planning requires the model to hold two simultaneous representations:
(a) the target topic (workweek economics) and (b) an abstract meta-schema of what evaluation
entails (consider\,$\to$\,analyze\,$\to$\,weigh\,$\to$\,verdict).
Secondary singular modes carry this meta-schema.
At $k=1$, only the topic is retained; the model cannot sustain the evaluative meta-structure
and generates the lower-rank fallback: open-ended questions that gesture toward evaluation
without performing it.

The $+7.7$ word-count inflation is a direct consequence: questions require more words
than imperative directives to convey the same evaluative intent.
Conversely, at $k=5$, step-structure recovery reaches \textbf{86\%}---the fastest convergence
of any level---because Evaluating tasks, once the full meta-schema is available,
converge to a narrow formulaic structure (weigh pros and cons, reach verdict).

\begin{figure}[h]
    \centering
    \begin{subfigure}{0.48\textwidth}
        \includegraphics[width=\textwidth]{/Users/vedantmaheshwari/.gemini/antigravity/scratch/neurIPS_BTL_multilayer/svd_btl_results/plots/logit_diff_vs_btl.png}
        \caption{Mean logit difference across BTL levels and $k$ values. The W-shape mirrors the text divergence pattern.}
    \end{subfigure}
    \begin{subfigure}{0.48\textwidth}
        \includegraphics[width=\textwidth]{/Users/vedantmaheshwari/.gemini/antigravity/scratch/neurIPS_BTL_multilayer/svd_btl_results/plots/dominance_vs_btl.png}
        \caption{Top-1 singular value dominance per layer and BTL level. Dominance is layer-intrinsic, not prompt-dependent.}
    \end{subfigure}
    \caption{Logit perturbation and spectral dominance. The W-shape in logit difference is consistent with the interrogative-collapse finding at Evaluating: evaluative prompts concentrate more
    attentional energy into $\sigma_1$ (high dominance), making them intrinsically more robust to $k=1$ truncation despite their stylistic volatility.}
\end{figure}

%% ── ── ── ── ── ── ── ── ── ── ── ── ── ── ── ── ── ── ──
\subsection{Creating: Contextual Genericization and Constraint Binding}

Creative synthesis prompts expose the most nuanced failure mode.
Consider the mountainous city transportation example:

\begin{quote}
\small
\textbf{Baseline:} Step 3: ``Research and identify suitable transportation modes
  for mountainous areas, \underline{such as cable cars or gondolas}.'' \\[4pt]
\textbf{$k=1$:}  Step 3: ``Determine the type of transportation system that would be
  most \underline{efficient and cost-effective}.'' \\[4pt]
\textbf{$k=3$:}  Step 3: ``Determine the type of transportation system
  (e.g., buses, trains, \underline{or cable cars}).'' \\[4pt]
\textbf{$k=5$:}  Step 3: ``Determine the type of transportation system most suitable
  \underline{for the mountainous terrain}.''
\end{quote}

The truncation trajectory reveals a \textbf{constraint binding hierarchy}:
\begin{enumerate}
    \item At $k=1$: The domain constraint (mountainous terrain) disappears entirely,
          replaced by generic optimisation language (efficient, cost-effective).
    \item At $k=3$: Domain-specific options (cable cars) partially re-emerge,
          but without the terrain justification.
    \item At $k=5$: The terrain constraint is restored, but without creative synthesis
          (no specific mention of cable cars or gondolas).
    \item At baseline: Domain constraint $\times$ transport engineering $\to$ novel synthesis
          (cable~cars, gondolas).
\end{enumerate}

This hierarchy demonstrates that cross-domain constraint binding---holding ``mountainous terrain''
and ``transport engineering'' in joint attention simultaneously to produce a contextualised
synthesis---requires progressively higher singular modes.
Each additional mode restores one layer of constraint binding specificity.
Creative tasks, by definition, require the binding of multiple domain constraints, explaining
why step-structure preservation is 0\% at $k=3$ and $k=4$---the worst of all levels.

%% ── ── ── ── ── ── ── ── ── ── ── ── ── ── ── ── ── ── ──
\subsection{The Five Generation Archetypes}

The six BTL levels reduce to five distinct generation archetypes under SVD truncation,
ordered by increasing rank requirement:

\begin{table}[h]
\centering
\caption{Generation archetypes under SVD truncation, from most to least rank-1 compatible.}
\label{tab:archetypes}
\begin{tabular}{lp{3.2cm}p{3.5cm}p{2.8cm}}
\toprule
\textbf{Archetype} & \textbf{BTL Levels} & \textbf{Effect at $k=1$} & \textbf{Recovery at $k=5$} \\
\midrule
Formulaic Immunity   & Applying (math)        & Perfect or near-perfect preservation & Perfect \\
Template Preservation & Remembering            & Qualifier compression; scaffold intact & Near-full \\
Lexical Substitution  & Analyzing              & Synonym swaps; logic chain preserved & Full \\
Framework Shift       & Understanding          & Register change (declarative→narrative) & Partial \\
Interrogative Collapse & Evaluating            & Imperative plan→exploratory questions; $+$7.7 words & 86\% step recovery \\
Contextual Genericization & Creating           & Domain constraints dropped; generic templates & Partial constraint restoration \\
\bottomrule
\end{tabular}
\end{table}

%% ── ── ── ── ── ── ── ── ── ── ── ── ── ── ── ── ── ── ──
\subsection{The Chaotic Trajectory Theorem: Spectral Rank Does Not Predict Text Divergence}

A critical control result bridges our spectral and text analyses.
We compute Pearson $r$ between Layer 31's effective rank and two text-divergence
measures across all 42 prompts:

\begin{align}
    r\bigl(r_\mathrm{eff}^{L31},\; \text{prefix match}\bigr) &= -0.106 \\
    r\bigl(r_\mathrm{eff}^{L31},\; D_\mathrm{KL}\bigr)       &= -0.006
\end{align}

The correlation is effectively zero.
A prompt with $r_\mathrm{eff}^{L31} = 4.6$ can diverge completely from baseline (0\% prefix);
a prompt with $r_\mathrm{eff}^{L31} = 5.5$ can match 41\% of baseline tokens.
The \emph{magnitude} of the initial spectral perturbation does not predict the
\emph{trajectory} of autoregressive divergence.

This is characteristic of a \textbf{chaotic dynamical system}: sensitive dependence on initial
conditions drives exponential divergence throughout the generation trajectory, but the divergence
is governed by the trajectory's Lyapunov structure, not by the norm of the initial perturbation.
SVD truncation acts as a small perturbation to the initial condition of the generation system;
the subsequent text trajectory compound-diverges through 200 autoregressive steps in a manner
that is orthogonal to the initial perturbation magnitude.

The practical implication is stark: \textbf{SVD truncation is not lossy compression of knowledge;
it is a stochastic perturbation of the generation style trajectory.}
The model's facts survive rank-1 reduction.
Only its stylistic routing---register, framing, constraint binding, meta-schema selection---is
displaced. This makes low-rank attention approximation viable for factual retrieval, classification,
and mathematical computation, but unreliable for tasks requiring sustained stylistic or
structural coherence.
